% Minimal preamble for Open Logic Project sources (pinned commit 1e960beff9ed7835bf3e3f1335e21af3439cd107).
%
% OLP's own master preamble cannot run on the project host: TeX Live 2025 there lacks
% prooftree.sty, ntheorem.sty, extarrows.sty, mathdots.sty, isomath.sty and subfiles.sty.
% Following the project's own rule ("if the author preamble needs packages the toolchain lacks,
% do not invent macros -- use a minimal working preamble"), this file declares, verbatim in
% spirit and numbering behaviour, exactly the theorem-like environments OLP defines in
% sty/open-logic-book-envs.sty (thm/lem/prop/cor/defn/ex/prob/rem/note/case/axiom/conv),
% plus the shorthand accents OLP uses (`!!a`, `!!b`, `!!{...}`).
%
% Author text is untouched; only the packaging is ours.
\documentclass[10pt,a4paper]{article}

\usepackage{amsmath}
\usepackage{amssymb}
\usepackage{amsthm}
\usepackage{graphicx}
\usepackage{hyperref}
\usepackage{enumitem}
\usepackage[framemethod=default]{mdframed}

% ---- OLP shorthand accents (from sty/open-logic.sty) ----
% NB: \! is taken by LaTeX (negative thin space), so we only define the !!a / !!b families
% and the \olup helper used by them.
\newcommand*{\olup}[1]{\MakeUppercase{#1}}
\def\!!!a#1{\olup{#1}}
\def\!!a#1{\olup{#1}}
\def\!!b#1{\olup{#1}}

% ---- theorem-like environments, same family numbering as OLP ----
\theoremstyle{plain}
\newtheorem{thm}{Theorem}[section]
\newtheorem{lem}[thm]{Lemma}
\newtheorem{prop}[thm]{Proposition}
\newtheorem{cor}[thm]{Corollary}
\newtheorem{axiom}{Axiom}

\theoremstyle{definition}
\newtheorem{defn}[thm]{Definition}
\newtheorem{ex}[thm]{Example}
\newtheorem{prob}{Problem}[section]
\newtheorem{probd}{Problem}
\newtheorem{conv}[thm]{Convention}

\theoremstyle{remark}
\newtheorem{rem}[thm]{Remark}
\newtheorem{note}[thm]{Note}
\newtheorem{case}[thm]{Case}

% ---- environments OLP defines as wrappers ----
\newenvironment{probtag}[1]{\begin{prob}}{\end{prob}}
\newenvironment{explain}{}{}
\newenvironment{explain*}{}{}
\newenvironment{editorial}{}{}
\newenvironment{digress}{}{}
\newenvironment{digress*}{}{}
\newenvironment{history}{}{}
\newenvironment{reference}{}{}
\newcommand*{\olfileid}[3]{}
\newcommand*{\olsection}[1]{\section*{#1}}
\newcommand*{\olchapter}[3]{}
\newcommand*{\olpart}[2]{}
\newcommand*{\olimport}[1]{}
\newcommand*{\olinclude}[1]{}
\newcommand*{\olpath}[1]{}
\newcommand*{\ollangid}[1]{}
\newcommand*{\problemtag}[1]{}
\newcommand*{\problemauthor}[1]{}
\newcommand*{\ollabel}[1]{\label{#1}}
% \olref takes 0-3 OPTIONAL arguments plus one mandatory label:
%     \olref{lab}  |  \olref[from]{lab}  |  \olref[from][to]{lab}  |  \olref[a][b][c]{lab}
% A single-argument definition silently swallows "[sth]" and emits a broken reference
% (measured on the host: the text came out as "??iso]x"). The chain below consumes the
% brackets and cross-references only the mandatory argument.
\makeatletter
\newcommand{\olref}{\@ifnextchar[{\olref@i}{\olref@final}}
\def\olref@i[#1]{\@ifnextchar[{\olref@ii}{\olref@final}}
\def\olref@ii[#1]{\@ifnextchar[{\olref@iii}{\olref@final}}
\def\olref@iii[#1]{\olref@final}
\def\olref@final#1{\ref{#1}}
\makeatother
\newcommand*{\olnewpage}{\clearpage}

% ===================== OLP compatibility layer (1) =====================

\makeatletter

% ---------------- proof trees (OLP's prooftree is a bussproofs derivative) ----------------
\providecommand{\AxiomC}[1]{}
\providecommand{\UnaryInfC}[1]{}
\providecommand{\BinaryInfC}[1]{}
\providecommand{\TrinaryInfC}[1]{}
\providecommand{\QuaternaryInfC}[1]{}
\providecommand{\QuinaryInfC}[1]{}
\providecommand{\RightLabel}[1]{}
\providecommand{\LeftLabel}[1]{}
\providecommand{\doubleLine}{}
\providecommand{\noLine}{}
\providecommand{\DisplayProof}{}
\providecommand{\fCenter}{}
\newenvironment{prooftree}{}{}
\newenvironment{ltprooftree}{}{}
\newenvironment{oltableau}{}{}

% ---------------- logical notation ----------------
\providecommand{\lif}{\mathbin{\rightarrow}}
\providecommand{\lforall}{\mathord{\forall}}
\providecommand{\lexists}{\mathord{\exists}}
\providecommand{\ltrue}{\top}
\providecommand{\lfalse}{\bot}
\providecommand{\Ltrue}{\top}
\providecommand{\Lfalse}{\bot}
\providecommand{\Sequent}{\mathrel{\Rightarrow}}
\providecommand{\Proves}{\mathrel{\vdash}}
\providecommand{\Entails}{\mathrel{\models}}
\providecommand{\Subst}[3]{#1[#2/#3]}
\providecommand{\sFmla}[2]{#1(#2)}
\providecommand{\pFmla}[2]{#1(#2)}
\providecommand{\formula}{\relax}
\providecommand{\True}{\top}
\providecommand{\False}{\bot}
\providecommand{\Args}[1]{#1}
\providecommand{\Ts}{\ensuremath{\top}}
\providecommand{\Fi}{\ensuremath{\bot}}
\providecommand{\lfalsum}{\bot}

% ---------------- objects, structures, metatheory ----------------
\providecommand{\Obj}{}
\providecommand{\Struct}{\ensuremath{\mathfrak{A}}}
\providecommand{\Domain}{\ensuremath{|\mathfrak{A}|}}
\providecommand{\Assign}{\ensuremath{s}}
\providecommand{\Sat}[2]{\ensuremath{#1 \models #2}}
\providecommand{\mSat}[3]{\ensuremath{#1 \models #2[#3]}}
\providecommand{\mModel}[2]{\ensuremath{#1 \models #2}}
\providecommand{\Lang}{\ensuremath{\mathcal{L}}}
\providecommand{\Th}{\ensuremath{\mathrm{Th}}}
\providecommand{\Log}{\ensuremath{\mathrm{L}}}
\providecommand{\tuple}[1]{\langle #1 \rangle}
\providecommand{\num}[1]{\underline{#1}}
\providecommand{\cfind}[1]{\varphi_{#1}}
\providecommand{\Nat}{\ensuremath{\mathbb{N}}}
\providecommand{\Mult}{\ensuremath{\mathrm{Mult}}}
\providecommand{\Add}{\ensuremath{\mathrm{Add}}}
\providecommand{\Exp}{\ensuremath{\mathrm{Exp}}}
\providecommand{\Zero}{\ensuremath{0}}
\providecommand{\Succ}{\ensuremath{\mathrm{S}}}
\providecommand{\lambd}{\ensuremath{\lambda}}
\providecommand{\bredpar}{\mathrel{\rightarrow_\beta}}
\providecommand{\bcd}[1]{\ensuremath{\ulcorner #1 \urcorner}}
\providecommand{\Set}[1]{\ensuremath{\{#1\}}}
\providecommand{\indcase}[1]{\par\noindent\emph{Case #1}.\quad}
\providecommand{\indfrm}[1]{\par\noindent\emph{Form #1}.\quad}
\providecommand{\iftag}[1]{#1}
\providecommand{\tagitem}[1]{\item[#1]}
\providecommand{\Elim}[1]{}
\providecommand{\Intro}[1]{}
\providecommand{\LeftR}[1]{}
\providecommand{\RightR}[1]{}
\providecommand{\TAss}[1]{}
\providecommand{\TRule}[1]{}

\makeatother
\providecommand{\olpath}{.}

% ===================== OLP compatibility layer (2) =====================
% --- second compatibility block: notation and shorthands used by the frozen OLP units ---
\providecommand{\CutCS}{\mathrm{Cut}}
\providecommand{\LeftSubproofLabel}{\relax}
\providecommand{\RightSubproofLabel}{\relax}
\providecommand{\maeh}{\relax}
\providecommand{\cheight}{\relax}
\providecommand{\ande}{\mathrm{and}_E}
\providecommand{\andi}{\mathrm{and}_I}
\providecommand{\ore}{\mathrm{or}_E}
\providecommand{\ori}{\mathrm{or}_I}
\providecommand{\typeof}{\mathord{:}}
\providecommand{\pValue}[1]{#1}
\providecommand{\Sent}{\mathrm{Sent}}
\providecommand{\cardplus}{\mathbin{+}}
\providecommand{\cardtimes}{\mathbin{\cdot}}
\providecommand{\cardexpo}{\mathbin{\hat{\ }}}
\providecommand{\cardnless}{\not<}
\providecommand{\cardle}{\leq}
\providecommand{\funfromto}[3]{#1\colon #2 \to #3}
\providecommand{\funimage}[1]{\operatorname{im}(#1)}
\providecommand{\fdefined}{\mathord{\downarrow}}
\providecommand{\canonord}{\operatorname{ord}}
\providecommand{\ordtimes}{\mathbin{\cdot}}
\providecommand{\ordeq}[2]{#1 = #2}
\providecommand{\ordneq}[2]{#1 \neq #2}
\providecommand{\ordtype}[1]{\operatorname{ord}(#1)}
\providecommand{\ordsucc}{\mathord{+1}}
\providecommand{\ordunion}{\cup}
\providecommand{\onesphere}{S^1}
\providecommand{\rotationsgroup}{\mathrm{SO}(2)}
\providecommand{\equivrep}[1]{[#1]}
\providecommand{\rlexless}{\prec}
\providecommand{\ZFC}{\mathrm{ZFC}}
\providecommand{\ZF}{\mathrm{ZF}}
\providecommand{\Rat}{\mathbb{Q}}
\providecommand{\Real}{\mathbb{R}}
\providecommand{\PosInt}{\mathbb{Z}^{+}}
\providecommand{\Id}{\mathrm{id}}
\providecommand{\Bin}{\mathrm{Bin}}
\providecommand{\pto}{\rightharpoonup}
\providecommand{\dom}{\operatorname{dom}}
\providecommand{\ran}{\operatorname{ran}}
\providecommand{\len}{\operatorname{len}}
\providecommand{\concat}{\mathbin{{}^\frown}}
\providecommand{\card}[1]{\lvert #1 \rvert}
\providecommand{\Pow}{\mathcal{P}}
\providecommand{\Setabs}[2]{\{\, #1 \mid #2 \,\}}
\providecommand{\comp}{\circ}
\providecommand{\Char}{\chi}
\providecommand{\Prf}{\operatorname{Prf}}
\providecommand{\Prov}{\operatorname{Prov}}
\providecommand{\Proj}{\operatorname{proj}}
\providecommand{\Gn}{\operatorname{gn}}
\providecommand{\gn}{\operatorname{gn}}
\providecommand{\Value}[1]{\operatorname{val}(#1)}
\providecommand{\Deduce}{\mathrm{D}}
\providecommand{\Dual}{\mathrm{d}}
\providecommand{\Ax}{\mathrm{Ax}}
\providecommand{\pSat}[3]{#1 \models #2\,[#3]}
\providecommand{\pAssign}{\sigma}
\providecommand{\bforall}{\forall}
\providecommand{\bexists}{\exists}
\providecommand{\liff}{\mathbin{\leftrightarrow}}
\providecommand{\FV}{\operatorname{FV}}
\providecommand{\UnaryInf}{\ensuremath{\mid}}
\providecommand{\BinaryInf}{\ensuremath{\mid}}
\providecommand{\depth}{\operatorname{depth}}
\providecommand{\oliflabeldef}[3]{#1}
\providecommand{\TMblank}{\sqcup}
\providecommand{\TMstroke}{\mid}
\providecommand{\TMstay}{\mathrm{S}}
\providecommand{\TMright}{\mathrm{R}}
\providecommand{\TMleft}{\mathrm{L}}
\providecommand{\TMendtape}{\triangleright}
\providecommand{\umin}{\mu}
\providecommand{\Sum}{\sum}
\providecommand{\citep}[1]{\cite{#1}}
\providecommand{\citet}[1]{\cite{#1}}
\providecommand{\citeauthor}[1]{\cite{#1}}
\providecommand{\citeyear}[1]{\cite{#1}}

% ===================== OLP compatibility layer (3) =====================
% --- third compatibility block: remaining notation seen in the frozen OLP units ---
\providecommand{\iso}{\cong}
\providecommand{\nszero}{\mathrm{0}}
\providecommand{\nssucc}{\mathrm{S}}
\providecommand{\nsplus}{\mathbin{+}}
\providecommand{\nssless}{<}
\providecommand{\nsless}{<}
\providecommand{\PL}{\mathrm{PL}}
\providecommand{\Nec}{\mathord{\Box}}
\providecommand{\RK}{\mathrm{RK}}
\providecommand{\ST}{\mathrm{ST}}
\providecommand{\SSubst}[3]{#1[#2/#3]}
\providecommand{\subst}[3]{#1[#2/#3]}
\providecommand{\tagfalse}{\mathrm{F}}
\providecommand{\cutr}{\mathrm{Cut}}
\providecommand{\cardneq}{\neq}
\providecommand{\issucc}[1]{\operatorname{succ}(#1)}
\providecommand{\tup}[1]{\langle #1 \rangle}
\providecommand{\Var}{\mathrm{Var}}
\providecommand{\Bew}{\operatorname{Bew}}
\providecommand{\Con}{\operatorname{Con}}
\providecommand{\Cons}{\operatorname{Cons}}
\providecommand{\pDeduce}{\mathrm{D}}
\providecommand{\psat}{\models}
\providecommand{\den}[1]{\llbracket #1 \rrbracket}
\providecommand{\Sem}[2]{\llbracket #1 \rrbracket_{#2}}
\providecommand{\DomainOf}[1]{\lvert #1 \rvert}
\providecommand{\Form}{\mathrm{Form}}
\providecommand{\Terms}{\mathrm{Term}}
\providecommand{\Atom}{\mathrm{Atom}}
\providecommand{\FVars}[1]{\operatorname{FV}(#1)}
\providecommand{\fv}[1]{\operatorname{fv}(#1)}
\providecommand{\tuplep}[1]{\langle #1 \rangle}
\providecommand{\restrict}[2]{#1 \upharpoonright #2}
\providecommand{\sto}{\rightarrow}
\providecommand{\FmlaSet}[1]{\mathrm{Fmla}_{#1}}
\providecommand{\Thm}{\mathrm{Thm}}
\providecommand{\Eval}{\operatorname{val}}
\providecommand{\blank}{\sqcup}
\providecommand{\blankTape}{\sqcup}
\providecommand{\Consis}{\operatorname{Con}}
\providecommand{\drs}{\mathrm{d}}
\providecommand{\VRule}{\mid}
\providecommand{\Bet}{\mathrm{B}}
\providecommand{\Seg}[1]{\mathrm{Seg}_{#1}}
\providecommand{\Enums}{\mathrm{En}}
\providecommand{\Z}{\mathbb{Z}}
\providecommand{\N}{\mathbb{N}}
\providecommand{\abs}[1]{\lvert #1 \rvert}
\providecommand{\floor}[1]{\lfloor #1 \rfloor}
\providecommand{\ceil}[1]{\lceil #1 \rceil}

% ===================== OLP compatibility layer (4) =====================
% --- fourth compatibility block: harvested from the per-candidate compile logs ---
\providecommand{\beredpar}{\mathrel{\rightarrow_\beta}}
\makeatletter
\def\bered{\@ifnextchar[{\bered@i}{\mathrel{\rightarrow_\beta}}}
\def\bered@i[#1]{\mathrel{\rightarrow_\beta}}
\def\bred{\@ifnextchar[{\bred@i}{\mathrel{\rightarrow_\beta}}}
\def\bred@i[#1]{\mathrel{\rightarrow_\beta}}
\makeatother
\providecommand{\beredone}{\mathrel{=_\beta}}
\providecommand{\bredone}{\mathrel{=_\beta}}
\providecommand{\xredone}{\mathrel{=_\beta}}
\providecommand{\aconvone}{\mathrel{=_\alpha}}
\providecommand{\cardless}{<}
\providecommand{\tagrefs}[1]{#1}
\providecommand{\ordplus}{\mathbin{+}}
\providecommand{\funrestrictionto}[2]{#1 \upharpoonright #2}
\providecommand{\fundefined}{\mathord{\uparrow}}
\providecommand{\Trm}{\mathrm{Trm}}
\providecommand{\OProv}{\operatorname{Prov}}
\providecommand{\OPrf}{\operatorname{Prf}}
\providecommand{\aeq}{\mathrel{=}}
\providecommand{\mClass}[1]{[#1]}
\providecommand{\equivclass}[1]{[#1]}
\providecommand{\pheight}{\operatorname{ht}}
\providecommand{\phy}{\operatorname{ht}}
\providecommand{\Cut}{\mathrm{Cut}}
\providecommand{\QuantRank}{\operatorname{qr}}
\providecommand{\elemequiv}{\mathrel{\equiv}}
\providecommand{\PIso}{\mathrm{PIso}}
\providecommand{\Theory}{\mathrm{Th}}
\providecommand{\supstrict}{\subsetneq}
\providecommand{\trcl}{\operatorname{trcl}}
\providecommand{\cref}[1]{\ref{#1}}
\providecommand{\Discharge}{\mathrm{Dis}}
\providecommand{\Weakening}{\mathrm{Wk}}
\providecommand{\ZFminus}{\mathrm{ZF}^{-}}
\providecommand{\RK}{\mathrm{RK}}
\makeatletter
\renewcommand{\Dual}{\@ifnextchar[{\Dual@i}{\ifmmode\mathrm{d}\else$\mathrm{d}$\fi}}
\def\Dual@i[#1]{\ifmmode\mathrm{d}[#1]\else$\mathrm{d}[#1]$\fi}
\makeatother
\providecommand{\RightR}{\mathrm{R}}
\providecommand{\LeftR}{\mathrm{L}}
\providecommand{\Intro}{\mathrm{I}}
\providecommand{\Axiom}{\mathrm{Ax}}
\providecommand{\formula}[1]{#1}
\providecommand{\stared}{\ast}
\providecommand{\oset}[2]{\{\, #1 \mid #2 \,\}}
\providecommand{\setof}[2]{\{\, #1 \mid #2 \,\}}
\providecommand{\pow}[1]{\mathcal{P}(#1)}
\providecommand{\ordpair}[2]{\langle #1, #2 \rangle}
\providecommand{\ordtriple}[3]{\langle #1, #2, #3 \rangle}
\providecommand{\union}{\cup}
\providecommand{\inter}{\cap}
\providecommand{\seq}[1]{\langle #1 \rangle}
\providecommand{\SuchThat}{\mid}
\providecommand{\where}{\mid}
\providecommand{\TVar}{\mathrm{TVar}}
\providecommand{\Bool}{\mathrm{Bool}}
\providecommand{\Val}{\operatorname{Val}}
\providecommand{\FV}{\operatorname{FV}}
\providecommand{\TL}[1]{\mathrm{TL}}
\providecommand{\xstared}{\ast}

% ===================== OLP compatibility layer (5) =====================
% --- fifth compatibility block: rescues the quarantined OLP candidates ---
\providecommand{\ident}{\operatorname{id}}
\providecommand{\setrank}{\operatorname{rank}}
\providecommand{\cardfont}{\mathfrak{c}}
\providecommand{\Complement}{\overline}
\providecommand{\QR}{\mathrm{QR}}
\providecommand{\becd}{\ensuremath{\ulcorner}}
\providecommand{\disjointsum}{\mathbin{\oplus}}
\providecommand{\cardeq}[2]{#1 = #2}
\providecommand{\xred}{\mathrel{\rightarrow}}
\providecommand{\cardless}{<}
\providecommand{\cardplus}{\mathbin{+}}
\providecommand{\cardtimes}{\mathbin{\cdot}}
\providecommand{\cardexpo}{\mathbin{\hat{\ }}}
\providecommand{\cardneq}{\neq}
\providecommand{\cardle}{\leq}
\providecommand{\cardnless}{\not<}
% --- sixth compatibility block: names reported by real compile logs ---
\providecommand{\ORefut}{\operatorname{Refut}}
\providecommand{\ORProv}{\operatorname{Prov}}
\providecommand{\cutrank}{\operatorname{cr}}
\providecommand{\Refut}{\operatorname{Refut}}
\providecommand{\DischargeRule}{\mathrm{Dis}}
\providecommand{\Contraction}{\mathrm{C}}
\providecommand{\bottomAlignProof}{\relax}
\providecommand{\FalseCl}{\bot}
\providecommand{\red}{\mathrel{\rightarrow}}
\providecommand{\varAssign}{\sigma}
\providecommand{\defiff}{\mathrel{\leftrightarrow}}
\providecommand{\FalseInt}{\bot}
\providecommand{\fn}{\operatorname{fn}}
\providecommand{\stared}{\ast}
\providecommand{\dotted}{\cdot}
\providecommand{\Gamma}{\Gamma}
% --- sixth compatibility block: names reported by real compile logs ---
\providecommand{\ORefut}{\operatorname{Refut}}
\providecommand{\ORProv}{\operatorname{Prov}}
\providecommand{\cutrank}{\operatorname{cr}}
\providecommand{\Refut}{\operatorname{Refut}}
\providecommand{\DischargeRule}{\mathrm{Dis}}
\providecommand{\Contraction}{\mathrm{C}}
\providecommand{\bottomAlignProof}{\relax}
\providecommand{\FalseCl}{\bot}
\providecommand{\red}{\mathrel{\rightarrow}}
\providecommand{\varAssign}{\sigma}
\providecommand{\defiff}{\mathrel{\leftrightarrow}}
\providecommand{\FalseInt}{\bot}
\providecommand{\fn}{\operatorname{fn}}
\providecommand{\stared}{\ast}
\providecommand{\dotted}{\cdot}
\providecommand{\Gamma}{\Gamma}
\begin{document}

% ==== OLP file: content/set-theory/z/arbintersections.tex ====


\olfileid{sth}{z}{arbintersections}
\olsection{Appendix: Closure, Comprehension, and Intersection}

In \olref[sth][z][milestone]{sec}, we suggested that you should look
back through the na\"ive work of \olref[sfr][][]{part} and check that
it can be carried out in~$\Zminus$. If you followed that advice,
\emph{one} point might have tripped you up: the use of
\emph{intersection} in Dedekind's treatment of \emph{closures}. 

Recall from \olref[sfr][infinite][dedekind]{Closure} that
\begin{align*}
  \closureofunder{f}{o} & = \bigcap\Setabs{X}{o \in X \text{ and $X$ is
$f$-closed}}.
\intertext{The general shape of this is a definition of the form:}
  C & = \bigcap\Setabs{X}{\phi(X)}.
\end{align*}
But this should ring alarm bells: since Na\"ive Comprehension fails,
there is no guarantee that $\Setabs{X}{\phi(X)}$ exists. It looks
dangerously, then, like such definitions are \emph{cheating}. 

Fortunately, they are not cheating; or rather, if they \emph{are}
cheating as they stand, then we can engage in some honest toil to
render them kosher. That honest toil was foreshadowed in
\olref[sth][z][sep]{prop:intersectionsexist}, when we explained why
$\bigcap A$ exists for any $A \neq \emptyset$. But we will spell it out
explicitly.

Given Extensionality, if we attempt to define $C$ as
$\bigcap\Setabs{X}{\phi(X)}$, all we are really asking is for an
object $C$ which obeys the following:
\begin{equation}\ollabel{bicondelimarbintersection}
	\forall x(x \in C \liff \forall X(\phi(X) \lif x \in X))\tag{*}
\end{equation}
Now, suppose there is \emph{some} set, $S$, such that $\phi(S)$. Then
to deliver \olref{bicondelimarbintersection}, we can simply define $C$
using \emph{Separation}, as follows:
\[
	C = \Setabs{x \in S}{\forall X(\phi(X) \lif x \in X)}.
\]
We leave it as an exercise to check that this definition yields
\olref{bicondelimarbintersection}, as desired. 
%  fix $x$. First suppose that $x \in C$, as (re)defined; then both $x
%  \in S$ and $\forall X(\phi(X) \lif x \in X)$; but since
%  $\phi(S)$, this reduces to the condition that $\forall X(\phi(X)
%  \lif x \in X)$. Conversely, suppose $\forall X(\phi(X) \lif
%  x \in X)$; then, since $\phi(S)$, we also have that $x \in S$; so
%  $x \in c$, as (re)defined. 
And this general strategy will allow us to circumvent any apparent use
of Na\"ive Comprehension in defining intersections. In the particular
case which got us started on this line of thought, namely that of
$\closureofunder{f}{o}$, here is how that would work. We began the
proof of \olref[sfr][infinite][dedekind]{closureproperties} by noting
that $o \in \ran{f}\cup\{o\}$ and that $\ran{f} \cup \{o\}$ is
$f$-closed. So, we can define what we want thus:
\[
	\closureofunder{f}{o} = \Setabs{x \in \ran{f} \cup \{o\}}{(\forall X \ni o)(X \text{ is $f$-closed} \lif x \in X)}.
\]

% ==== OLP file: content/set-theory/z/infinity-again.tex ====


\olfileid{sth}{z}{infinity-again}	
\olsection{Infinity}

We already have enough axioms to ensure that there are infinitely many
sets (if there are any). For suppose some set exists, and so
$\emptyset$ exists (by \olref[sth][z][sep]{prop:emptyexists}). Now for
any set~$x$, the set $x \cup \{x\}$ exists by
\olref[sth][z][pairs]{prop:pairsconsequences}. So, applying this a few
times, we will get sets as follows:
\begin{enumerate}
	\item[0.] $\emptyset$ %& $0$ & stage $0$\\
	\item[1.] $ \{\emptyset\}$ %& $1$ & stage $1$\\
	\item[2.] $\{\emptyset, \{\emptyset\}\}$ %& $2$ & stage $2$\\
	\item[3.] $\{\emptyset, \{\emptyset\}, \{\emptyset, \{\emptyset\}\}\}$ %&$3$ & stage $3$\\
	\item[4.] $\{\emptyset, \{\emptyset\}, \{\emptyset, \{\emptyset\}\}, \{\emptyset, \{\emptyset\}, \{\emptyset, \{\emptyset\}\}\}\}$% & $4$ & stage $4$ 
\end{enumerate}
and we can check that each of these sets is distinct. 

We have started the numbering from $0$, for a few reasons. But one of
them is this. It is not that hard to check that the set we have
labelled ``$n$'' has exactly $n$ members, and (intuitively) is formed
at the $n$th stage. 

But. This gives us \emph{infinitely many} sets, but it does not
guarantee that there is an \emph{infinite set}, i.e., a set with
infinitely many members. And this really matters: unless we can find a
(Dedekind) infinite set, we cannot construct a Dedekind algebra. But
we want a Dedekind algebra, so that we can treat it as the set of
natural numbers. (Compare
\olref[sfr][infinite][dedekindsproof]{sec}.)

Importantly, the axioms we have laid down so far do \emph{not}
guarantee the existence of any infinite set. So we have to lay down a
new axiom:

\begin{axiom}[Infinity]
There is a set, $I$, such that $\emptyset \in I$ and $x \cup \{x\} \in I$ whenever $x \in I$.
\begin{align*}
	\exists I( & (\exists o \in I)\forall x\ x \notin o \land {}\\
	& (\forall x \in I)(\exists s \in I)\forall z(z \in s \liff (z \in x \lor z = x)))
\end{align*}
\end{axiom}

It is easy to see that the set $I$ given to us by the Axiom of
Infinity is Dedekind infinite. Its distinguished element is
$\emptyset$, and the injection on $I$ is given by $s(x) = x\cup
\{x\}$. Now,
\olref[sfr][infinite][dedekind]{thm:DedekindInfiniteAlgebra} showed
how to extract a Dedekind Algebra from a Dedekind infinite set; and we
will treat this as our set of natural numbers. More precisely:
\begin{defn}\ollabel{defnomega}
Let $I$ be any set given to us by the Axiom of Infinity. Let $s$
be the function $s(x) = x \cup \{x\}$. Let $\omega =
\closureofunder{s}{\emptyset}$. We call the members of $\omega$ the
\emph{natural numbers}, and say that $n$ is the result of $n$-many
applications of $s$ to $\emptyset$.
\end{defn}

You can now look back and check that the set labelled ``$n$'', a few
paragraphs earlier, will be treated \emph{as} the number $n$. 

We will discuss this significance of this stipulation in
\olref[sth][z][nat]{sec}. For now, it enables us to prove an
intuitive result:

\begin{prop}\ollabel{naturalnumbersarentinfinite}
No natural number is Dedekind infinite.
\end{prop}

\begin{proof}
The proof is by induction, i.e.,
\olref[sfr][infinite][induction]{thm:dedinfiniteinduction}. Clearly $0
= \emptyset$ is not Dedekind infinite. For the induction step, we will
establish the contrapositive: if (absurdly) $s(n)$ is Dedekind
infinite, then $n$ is Dedekind infinite. 

So suppose that $s(n)$ is Dedekind infinite, i.e., there is some
!!{injection} $f$ with $\ran{f}\subsetneq \dom{f} = s(n) = n \cup
\{n\}$. There are two cases to consider. 

\emph{Case 1: $n \notin \ran{f}$.} So $\ran{f} \subseteq n$, and $f(n) \in n$. Let $g = \funrestrictionto{f}{n}$; now $\ran{g} =
\ran{f} \setminus \{f(n)\} \subsetneq n = \dom{g}$. Hence $n$ is
Dedekind infinite. 

\emph{Case 2: $n \in \ran{f}$.} Fix $m \in \dom{f} \setminus \ran{f}$, and define a function $h$ with domain $s(n) = n \cup \{n\}$:
\[
h(x) = 
	\begin{cases}
		f(x) & \text{if }f(x) \neq n\\
		m & \text{if }f(x)=n
	\end{cases}
\]
So $h$ and $f$ agree everywhere, except that $h(f^{-1}(n)) = m
\neq n = f(f^{-1}(n))$. Since $f$ is !!a{injection}, $n \notin
\ran{h}$; and $\ran{h} \subsetneq \dom{h} = s(n)$. Now $n$ is
Dedekind infinite, using the argument of Case 1.
\end{proof}

The question remains, though, of how we might \emph{justify} the Axiom
of Infinity. The short answer is that we will need to add another
principle to the story we have been telling. That principle is as
follows:
\begin{enumerate}
	\item[] \stagesinf. There is an infinite stage. That is, there is a stage which (a) is not the first stage, and which (b) has some stages before it, but which (c) has no immediate predecessor.
\end{enumerate}
The Axiom of Infinity follows straightforwardly from this principle.
We know that natural number $n$ is formed at stage $n$. So the set
$\omega$ is formed at the first infinite stage. And $\omega$ itself
witnesses the Axiom of Infinity. 

This, however, simply pushes us back to the question of how we might
justify \stagesinf. As with \stagessucc, it was not an explicit part
of the story we told about the cumulative-iterative hierarchy. But
more than that: nothing in the very idea of an iterative hierarchy, in
which sets are formed stage by stage, forces us to think that the
process involves an \emph{infinite} stage. It seems perfectly coherent
to think that the stages are ordered like the natural numbers. 

This, however, gives rise to an obvious problem. In
\olref[sfr][infinite][dedekindsproof]{sec}, we considered
Dedekind's ``proof'' that there is a Dedekind infinite set (of
thoughts). This may not have struck you as very satisfying. But if
\stagesinf{} is not ``forced upon us'' by the iterative conception of
set (or by ``the laws of thought''), then we are still left without an
intrinsic justification for the claim that there is a Dedekind
infinite set.

There is much more to say here, of course. But hopefully you are now
at a point to start thinking about what it might \emph{take} to
justify an axiom (or principle). In what follows we will simply take
\stagesinf{} for granted.

% ==== OLP file: content/set-theory/z/milestone.tex ====


\olfileid{sth}{z}{milestone}
\olsection{$\Z^-$: a Milestone}

We will revisit \stagesinf{} in the next section. However, with the
Axiom of Infinity, we have reached an important milestone. We now have
all the axioms required for the theory $\Zminus$. In detail:

\begin{defn}
The theory $\Zminus$ has these axioms: Extensionality, Union, Pairs, Powersets, Infinity, and all instances of the Separation scheme.
\end{defn}

The name stands for \emph{Zermelo} set theory (\emph{minus} something
which we will come to later). Zermelo deserves the honour, since he
essentially formulated this theory in his
\citeyear{Zermelo1908Untersuchungen}.\footnote{For interesting
comments on the history and technicalities, see \citet[Appendix
A]{Potter2004}.}

This theory is powerful enough to allow us to do an enormous amount of
mathematics. In particular, you \emph{should} look back through
\olref[sfr][][]{part}, and convince yourself that everything we did,
na\"ively, could be done more formally within~$\Zminus$. (Once you
have done that for a bit, you might want to skip ahead and read
\olref[sth][z][arbintersections]{sec}.) So, henceforth,
and without any further comment, we will take ourselves to be working
in $\Zminus$ (at least).

% ==== OLP file: content/set-theory/z/nat.tex ====

	
\olfileid{sth}{z}{nat}
\olsection{Selecting our Natural Numbers}

%Let us take it for granted that the informal principle \stagesinf{}
%can be justified. It is also, though, worth commenting on the
%particular axiom we extracted from that informal principle. 

In \olref[infinity-again]{defnomega}, we explicitly defined
the expression ``natural numbers''. How should you understand this
stipulation? It is not a metaphysical claim, but just a decision to
\emph{treat} certain sets as the natural numbers. We touched upon reasons for thinking this in
\olref[sfr][rel][ref]{sec}, \olref[sfr][arith][ref]{sec} and
\olref[sfr][infinite][dedekindsproof]{sec}. But we can make these
reasons even more pointed.

Our Axiom of Infinity follows \citet{VonNeumann1925}. But here is
another axiom, which we could have adopted instead:

\begin{defish}
\emph{Zermelo's \citeyear{Zermelo1908Untersuchungen} Axiom of
Infinity.} There is a set $A$ such that $\emptyset \in A$ and
$(\forall x \in A)\{x\} \in A$. 
\end{defish}

Had we used Zermelo's axiom, instead of our (von Neumann-inspired)
Axiom of Infinity, we would equally well have been given a Dedekind
infinite set, and so a Dedekind algebra. On Zermelo's approach, the
distinguished element of our algebra would again have been $\emptyset$
(our surrogate for $0$), but the injection would have been given by
the map $x \mapsto \{x\}$, rather than $x \mapsto x \cup \{x\}$. The
simplest upshot of this is that Zermelo treats $2$ as
$\{\{\emptyset\}\}$, whereas we (with von Neumann) treat $2$ as
$\{\emptyset, \{\emptyset\}\}$. 

Why choose one axiom of Infinity rather than the other? The main
practical reason is that von Neumann's approach ``scales up'' to
handle transfinite numbers rather well. We will explore this from
\olref[ordinals][]{chap} onwards. However, from the simple
perspective of \emph{doing arithmetic}, both approaches would do
equally well. So if someone tells you that the natural numbers
\emph{are} sets, the obvious question is: \emph{Which sets are they?} 

This precise question was made famous by \citet{Benacerraf1965}. But
it is worth emphasising that it is just the most famous example of a
phenomenon that we have encountered many times already. The basic
point is this. Set theory gives us a way to \emph{simulate} a bunch of
``intuitive'' kinds of entities: the reals, rationals, integers, and
naturals, yes; but also ordered pairs, functions, and relations.
However, set theory never provides us with a \emph{unique} choice of
simulation. There are \emph{always} alternatives
which---straightforwardly---would have served us just as well. 

% ==== OLP file: content/set-theory/z/pairs.tex ====


\olfileid{sth}{z}{pairs}
\olsection{Pairs}

The next axiom to consider is the following:

\begin{axiom}[Pairs]
For any sets $a, b$, the set $\{a, b\}$ exists.
\[
	\forall a \forall b \exists P \forall x (x \in P \liff (x = a \lor x = b))
\]
\end{axiom}

Here is how to justify this axiom, using the iterative conception. Suppose $a$ is available at stage $S$, and $b$ is available at stage $T$. Let $M$ be whichever of stages $S$ and $T$ comes later. Then since $a$ and $b$ are both available at stage $M$, the set $\{a,b\}$ is a possible collection available at any stage after $M$ (whichever is the greater).

But hold on!{} Why assume that there \emph{are} any stages after $M$? If there are none, then our justification will fail. So, to justify Pairs, we will have to add another principle to the story we told in \olref[sth][z][story]{sec}, namely:
\begin{enumerate}
	\item[] \stagessucc. There is no last stage.
\end{enumerate}
Is this principle justified? Nothing in Shoenfield's story stated
\emph{explicitly} that there is no last stage. Still, even if it is
(strictly speaking) an extra addition to our story, it fits well with
the basic idea that sets are formed in stages. We will simply accept
it in what follows. And so, we will accept the Axiom of Pairs too.

Armed with this new Axiom, we can prove the existence of plenty more sets. For example:

\begin{prop}\ollabel{prop:pairsconsequences}
For any sets $a$ and $b$, the following sets exist:
	\begin{enumerate}
		\item\ollabel{singleton} $\{a\}$
		\item\ollabel{binunion} $a \cup b$
		\item\ollabel{tuples} $\tuple{a, b}$
	\end{enumerate}
\end{prop}

\begin{proof}
\olref{singleton}. By Pairs, $\{a, a\}$ exists, which is $\{a\}$ by
Extensionality.

\olref{binunion}. By Pairs, $\{a, b\}$ exists. Now $a \cup b = \bigcup
\{a, b\}$ exists by Union.

\olref{tuples}. By \olref{singleton}, $\{a\}$ exists. By Pairs, $\{a,
b\}$ exists. Now $\{\{a\}, \{a, b\}\} = \tuple{a, b}$ exists, by Pairs
again.
\end{proof}

\begin{prob}
Show that, for any sets $a, b, c$, the set $\{a, b, c\}$ exists.
\end{prob}

\begin{prob}
Show that, for any sets $a_1, \ldots, a_n$, the set $\{a_1, \ldots,
a_n\}$ exists.
\end{prob}

%\begin{proof} By two applications of Pairs, $\{\{a_1, a_2\}, \{a_1,
%   a_3\}\}$ exists. By Union and Extensionality, $\bigcup \{\{a_1,
%   a_2\}, \{a_1, a_3\}\} = \{a_1, a_2, a_3\}$ exists. Repeat this
%   trick as often as necessary. \end{proof}


% ==== OLP file: content/set-theory/z/powerset.tex ====


\olfileid{sth}{z}{power}
\olsection{Powersets}

We will proceed with another axiom:

\begin{axiom}[Powersets]
For any set $A$, the set $\Pow{A} = \Setabs{x}{x \subseteq A}$ exists.
\[
	\forall A \exists P \forall x(x \in P \liff (\forall z \in x)z \in A)
\]
\end{axiom}

Our justification for this is pretty straightforward. Suppose $A$ is
formed at stage $S$. Then all of $A$'s members were available before
$S$ (by \stagesacc). So, reasoning as in our justification for
Separation, every subset of $A$ is formed by stage $S$. So they are
all available, to be formed into a single set, at any stage after $S$.
And we know that there is some such stage, since $S$ is not the last
stage (by \stagessucc). So $\Pow{A}$ exists.

Here is a nice consequence of Powersets:

\begin{prop}\label{thm:Products}
Given any sets $A, B$, their Cartesian product $A \times B$ exists.
\end{prop}

\begin{proof}
The set $\Pow{\Pow{A \cup B}}$ exists by Powersets and
\olref[sth][z][pairs]{prop:pairsconsequences}. So by Separation, this
set exists:
\[
	C = \Setabs{z \in \Pow{\Pow{A \cup B}}}{(\exists x \in A)(\exists y \in B) z = \tuple{x, y}}.
\]
Now, for any $x \in A$ and $y \in B$, the set $\tuple{x, y}$ exists by
\olref[sth][z][pairs]{prop:pairsconsequences}. Moreover, since $x, y
\in A \cup B$, we have that $\{x\}, \{x, y\} \in \Pow{A \cup B}$, and
$\tuple{x,y} \in \Pow{\Pow{A \cup B}}$. So $A \times B = C$.
\end{proof}

In this proof, Powerset interacts with Separation. And that is no
surprise. Without Separation, Powersets wouldn't be a very
\emph{powerful} principle. After all, Separation tells us which
subsets of a set exist, and hence determines just how ``fat'' each
Powerset is.

\begin{prob}
Show that, for any sets $A, B$: (i) the set of all relations with
domain $A$ and range $B$ exists; and (ii) the set of all functions
from $A$ to $B$ exists.
\end{prob}

\begin{prob}
Let $A$ be a set, and let $\sim$ be an equivalence relation on $A$.
Prove that the set of equivalence classes under $\sim$ on $A$, i.e.,
$\equivclass{A}{\sim}$, exists.
\end{prob}


% ==== OLP file: content/set-theory/z/separation.tex ====


\olfileid{sth}{z}{sep}
\olsection{Separation}

We start with a principle to replace Na\"{i}ve Comprehension:

\begin{axiom}[Scheme of Separation] For every formula $\phi(x)$, this is an
axiom: for any $A$, the set $\Setabs{x \in A}{\phi(x)}$ exists.
\end{axiom}

Note that this is not a single axiom. It is a \emph{scheme} of axioms.
There are \emph{infinitely many} Separation axioms; one for every
formula $\phi(x)$. The scheme can equally well be (and normally is)
written down as follows:

\begin{defish}
For any formula $\phi(x)$ which does not contain ``$S$'', this is an
axiom:
\[
	\forall A \exists S \forall x(x \in S \liff (\phi(x) \land x \in A)).
\]
\end{defish}

In keeping with the convention noted at the start of
\olref[sth][][]{part}, the formulas~$\phi$ in the Separation axioms
may have parameters.\footnote{For an explanation of what this means,
see the discussion immediately after
\olref[sfr][infinite][induction]{natinductionschema}.}

Separation is immediately justified by our cumulative-iterative
conception of sets we have been telling. To see why, let $A$ be a set.
So $A$ is formed by some stage~$S$ (by \stageshier). Since $A$ was
formed at stage~$S$, all of $A$'s members were formed before stage $S$
(by \stagesacc). Now in particular, consider all the sets which are
members of $A$ and which also satisfy $\phi$; clearly all of these
sets, too, were formed before stage~$S$. So they are formed into a set
$\Setabs{x \in A}{\phi(x)}$ at stage~$S$ too (by \stagesacc).

Unlike Na\"ive Comprehension, this avoid Russell's Paradox. For we
cannot simply assert the existence of the set $\Setabs{x}{x \notin
x}$. Rather, \emph{given} some set~$A$, we can assert the existence of
the set $R_A = \Setabs{x \in A}{x \notin x}$. But all this proves is
that $R_A \notin R_A$ and $R_A \notin A$, none of which is very
worrying.

However, Separation has an immediate and striking consequence:

\begin{thm}\ollabel{thm:NoUniversalSet}
There is no \emph{universal} set, i.e., $\Setabs{x}{x = x}$ does not exist.
\end{thm}

\begin{proof}
For reductio, suppose $V$ is a universal set. Then by Separation, $R =
\Setabs{x \in V}{x \notin x} = \Setabs{x}{x \notin x}$ exists,
contradicting Russell's Paradox.
\end{proof}

The absence of a universal set---indeed, the open-endedness of the
hierarchy of sets---is one of the most fundamental ideas behind the
cumulative-iterative conception. So it is worth seeing that,
intuitively, we could reach it via a different route. A universal set
must be !!a{element} of itself. But, on our cumulative-iterative
conception, every set appears (for the first time) in the hierarchy at
the first stage immediately after all of its !!{element}s. But this
entails that \emph{no} set is self-membered. For any self-membered set
would have to first occur immediately after the stage at which it
first occurred, which is absurd. (We will see in
\olref[sth][spine][rank]{defnsetrank} how to make this explanation
more rigorous, by using the notion of the ``rank'' of a set. However,
we will need to have a few more axioms in place to do this.)

Here are a few more consequences of Separation and Extensionality.

\begin{prop}\ollabel{prop:emptyexists}
If any set exists, then $\emptyset$ exists.
\end{prop}

\begin{proof}
If $A$ is a set, $\emptyset = \Setabs{x \in A}{x \neq x}$ exists by Separation.
\end{proof}

\begin{prop}
$A \setminus B$ exists for any sets $A$ and $B$
\end{prop}

\begin{proof}
$A \setminus B = \Setabs{x \in A}{x \notin B}$ exists by Separation.
\end{proof}

It also turns out that (almost) arbitrary intersections exist:

\begin{prop}\ollabel{prop:intersectionsexist}
If $A \neq \emptyset$, then $\bigcap A = \Setabs{x}{(\forall y \in A)x \in y}$ exists.
\end{prop}

\begin{proof}
Let $A \neq \emptyset$, so there is some $c \in A$. Then $\bigcap A =
\Setabs{x}{(\forall y \in A)x \in y} = \Setabs{x \in c}{(\forall y \in
A)x \in y}$, which exists by Separation.
\end{proof}

Note the condition that $A \neq \emptyset$, though; for $\bigcap
\emptyset$ would be the universal set, vacuously, contradicting
\olref{thm:NoUniversalSet}.


% ==== OLP file: content/set-theory/z/story.tex ====


\olfileid{sth}{z}{story}
\olsection{The Story in More Detail}
	
In \olref[story][approach]{sec}, we quoted Schoenfield's
description of the process of set-formation. We now want to write down
a few more principles, to make this story a bit more precise. Here
they are:
\begin{enumerate}
\item[] \stageshier. Every set is formed at some stage.
\item[] \stagesord. Stages are ordered: some come \emph{before}
others.\footnote{We will actually assume---tacitly---that the stages
are \emph{well-ordered}. What this amounts to is explained in
\olref[ordinals][]{chap}. This is a substantial assumption. In fact,
using a very clever technique due to \citet{Scott1974}, this
assumption can be \emph{avoided} and then \emph{derived}. (This will
also explain why we should think that there is an initial stage.) We cannot go into that here; for more, see \citet{ButtonLT1}.} 
\item[] \stagesacc. For any stage $S$, and for any sets which were
formed \emph{before} stage~$S$: a set is formed at stage~$S$ whose
members are exactly those sets. Nothing else is formed at stage~$S$.
\end{enumerate}
These are informal principles, but we will be able to use them to
vindicate several of the axioms of Zermelo's set theory. 

(We should offer a word of caution. Although we will be presenting some
completely standard axioms, with completely standard names, the
italicized principles we have just presented have no particular names
in the literature. We simply monikers which we hope are helpful.)

% ==== OLP file: content/set-theory/z/union.tex ====


\olfileid{sth}{z}{union}
\olsection{Union}

\olref[sth][z][sep]{prop:intersectionsexist} gave us intersections.
But if we want arbitrary unions to exist, we need to lay down another
axiom:

\begin{axiom}[Union] For any set $A$, the set $\bigcup A =
\Setabs{x}{(\exists b \in A) x \in b}$ exists.
\[
	\forall A \exists U \forall x(x \in U \liff (\exists b \in A)x \in b)
\]
\end{axiom}

This axiom is also justified by the cumulative-iterative conception.
Let $A$ be a set, so $A$ is formed at some stage $S$ (by \stageshier).
Every member of $A$ was formed \emph{before}  $S$ (by \stagesacc); so,
reasoning similarly, every member of every member of $A$ was formed
before $S$. Thus all of \emph{those} sets are available before $S$, to
be formed into a set at $S$. And that set is just $\bigcup A$.


% ==== OLP file: content/set-theory/z/z.tex ====
% Part: set-theory
% Chapter: z



\olchapter{sth}{z}{Steps towards $\Z$}

\olimport{story}
\olimport{separation}
\olimport{union}
\olimport{pairs}
\olimport{powerset}
\olimport{infinity-again}
\olimport{milestone}
\olimport{nat}
\olimport{arbintersections}

\OLEndChapterHook

